%% ============================================================
%%  SECCIÓN 11 — Análisis de Gráficas
%%  Responsable: Vera Vivanco, Jesús Rodolfo
%% ============================================================

\section{Análisis de Gráficas}
\label{sec:analisis_graficas}

\subsection{Comparativa de los coeficientes \texorpdfstring{$\alpha$}{alpha} experimentales frente a los referentes teóricos --- Fig.~\ref{fig:grafica_alfa}}

Al observar la Fig.~\ref{fig:grafica_alfa}, se aprecia una distribución discreta de los coeficientes de dilatación lineal obtenidos para el cobre, el aluminio y el vidrio borosilicato. Los puntos experimentales siguen la tendencia de respuesta térmica esperada, mostrando de forma inequívoca que el aluminio posee la mayor susceptibilidad a la dilatación térmica, seguido del cobre y finalmente el vidrio borosilicato.

En cuanto a la posición de los puntos experimentales (puntos azules) con respecto a sus respectivas referencias bibliográficas (cruces rojas), se observan discrepancias particulares para cada muestra:
\begin{itemize}
    \item \textbf{Cobre:} El valor experimental de $\alpha_{\text{Cu}}$ subestima el valor bibliográfico teórico en un 5,8\%. La cruz roja se ubica ligeramente por encima del límite superior de la barra de error experimental. Esto sugiere la presencia de un error sistemático que sesga la medición por debajo de su valor esperado, el cual no llega a ser completamente compensado por la incertidumbre puramente instrumental.
    \item \textbf{Aluminio:} El valor experimental de $\alpha_{\text{Al}}$ sobreestima el valor de referencia en un 6,8\%. En este caso, la referencia teórica se encuentra justo por debajo de la barra de error, lo que indica un sesgo sistemático por exceso en la medición.
    \item \textbf{Vidrio borosilicato:} El punto experimental de $\alpha_{\text{vidrio}}$ muestra una concordancia excelente con el valor teórico de referencia ($E_r = 1{,}25\%$), con el valor de referencia ubicado cómodamente dentro de la barra de error.
\end{itemize}

La magnitud de las barras de error es consistente con la precisión de los instrumentos analizados en la Sección~\ref{sec:analisis_exp}. Específicamente, la barra de error del vidrio borosilicato es más ancha debido al peso relativo de la incertidumbre del transportador graduado ante un ángulo de giro pequeño ($\Delta\theta = 12,0^\circ \pm 0,5^\circ$). Para el cobre y el aluminio, el tamaño de las barras de error está dominado por la tolerancia del calibrador vernier al medir el diámetro milimétrico de la aguja. Las causas físicas detrás de las desviaciones sistemáticas observadas, tales como el deslizamiento mecánico de la aguja o la disipación térmica por convección en las paredes de los tubos, se analizan en detalle en la Sección~\ref{sec:conclusiones}. 